\section{ماژول kronecker\_tt\_lora (Kronecker-TT-LoRA Hybrid Compressor)}
\label{sec:module_kronecker_tt_lora}

\subsection{معرفی و انگیزه}
این مکانیزم بر مبنای مبانی نظری پیشرفته در ادبیات توجه و فشرده‌سازی \cite{ref_kroneckerttlora_2024} توسعه یافته است.

افزایش مقیاس مدل‌های زبانی ترنسفورمر نیازمند روش‌های تنظیم دقیق با بهره‌وری پارامتری بالا (\lr{PEFT}) و فشرده‌سازی هم‌زمان است. در حالی که روش \lr{LoRA} از یک فاکتورگیری ساده دو ماتریسی استفاده می‌کند، ظرفیت بیانی آن برای مدل‌سازی رفتارهای فرکانس بالا و تغییرات موضعی تانسورهای وزن کافی نیست. از طرف دیگر، شبکه‌های قطار تانسوری (\lr{Tensor-Train - TT}) توانایی بازنمایی ساختارهای تانسوری چندبعدی را با پارامترهای بسیار کم دارند اما در تطبیق‌پذیری سریع با داده‌های جدید دچار ایستایی گرادیان می‌شوند.

ماژول پیوندی \lr{kronecker\_tt\_lora} به عنوان یک نوآوری هیبریدی، حاصل‌ضرب هادامارد بین یک ساختار ضرب کرونکر و قطار تانسوری ۳-هسته‌ای را با یک شاخه کلاسیک کم‌رتبه \lr{LoRA} تلفیق می‌کند. این معماری سه مکانیزم متمایز را به صورت هم‌افزا به کار می‌گیرد: فاکتورهای کرونکر ساختار متناوب و سراسری وزن را ثبت می‌کنند، شبکه \lr{TT} همبستگی‌های غیرخطی چندوجهی را مدل‌سازی می‌نماید و شاخه \lr{LoRA} وظیفه انطباق سریع گرادیانی با وظیفه هدف را عهده‌دار می‌شود.

\subsection{مبانی نظری و فرمول‌بندی ریاضی}
در این ماژول، تغییرات یا مقدار مؤثر وزن لایه متراکم به صورت ترکیبی خطی از دو مؤلفه ساختاری تانسوری و انطباقی مدل‌سازی می‌گردد:
\begin{equation}
\Delta \mathbf{W} = \gamma \left( (\mathbf{A}_1 \otimes \mathbf{A}_2) \odot \mathcal{TT}(\mathbf{G}_1, \mathbf{G}_2, \mathbf{G}_3) \right)_{[:d_{\text{out}}, :d_{\text{in}}]} + \frac{\alpha}{r} (\mathbf{B}_{\text{lora}} \mathbf{A}_{\text{lora}})
\end{equation}
که در آن:
\begin{itemize}
    \item مؤلفه کرونکر حاصل ضرب دو ماتریس ارتوگونال با مقیاس طیفی $\mathbf{A}_1 \in \mathbb{R}^{m_1 \times n_1}$ و $\mathbf{A}_2 \in \mathbb{R}^{m_2 \times n_2}$ است به گونه‌ای که $m_1 m_2 \ge d_{\text{out}}$ و $n_1 n_2 \ge d_{\text{in}}$.
    \item عملگر تانسوری $\mathcal{TT}$ از سه هسته تانسوری $\mathbf{G}_1 \in \mathbb{R}^{1 \times d_1 \times r_1}$، $\mathbf{G}_2 \in \mathbb{R}^{r_1 \times d_2 \times r_2}$ و $\mathbf{G}_3 \in \mathbb{R}^{r_2 \times d_3 \times 1}$ با رتبه‌های داخلی $(1, r_1, r_2, 1)$ تشکیل شده است که ضرب داخلی آن‌ها تانسوری با ابعاد کلی $d_1 \times d_2 \times d_3 = (m_1 m_2) \cdot (n_1 n_2)$ ایجاد می‌نماید:
    \begin{equation}
    \mathcal{TT}_{i, j, k} = \sum_{a=1}^{r_1} \sum_{b=1}^{r_2} G_{1}(1, i, a) \, G_{2}(a, j, b) \, G_{3}(b, k, 1)
    \end{equation}
    \item عملگر $\odot$ نشان‌دهنده ضرب مؤلفه‌ای هادامارد است که اثر مدولاسیون ساختاری ایجاد می‌کند.
    \item شاخه آداپتور با $\mathbf{A}_{\text{lora}} \in \mathbb{R}^{r \times d_{\text{in}}}$ و $\mathbf{B}_{\text{lora}} \in \mathbb{R}^{d_{\text{out}} \times r}$ با ضریب مقیاس‌بندی $\frac{\alpha}{r}$ به خروجی افزوده می‌شود.
\end{itemize}

در نتیجه، وزن کل عملگر پیشرو به صورت زیر بر بردار ورودی $\mathbf{x}$ اعمال می‌گردد:
\begin{equation}
\mathbf{y} = \mathbf{x} \mathbf{W}_{\text{base}}^\top + \mathbf{x} \Delta \mathbf{W}^\top + \mathbf{b}
\end{equation}

\paragraph{تحلیل طیفی و نامساوی‌های پایداری خطا}
برای بررسی میزان دقت تقریب ماتریس وزن یا نگاشت توجه در ماژول \lr{kronecker\_tt\_lora}، تفاضل ماتریس فاکتورگیری‌شده از نگاشت متراکم اصلی بر اساس نامساوی هافدینگ-فون‌نیومان ارزیابی می‌گردد. فرض کنید $\mathbf{W} \in \mathbb{R}^{d_{\text{in}} \times d_{\text{out}}}$ ماتریس وزن هدف باشد که توسط تقریب کم‌مرتبه یا پراکنده با رتبه موثر $r$ جایگزین شده است. کران بالای خطای بازسازی بر مبنای نرم فروبنیوس به صورت زیر تصریح می‌شود:
\begin{equation}
\label{eq:spectral_bound_kronecker_tt_lora}
\| \mathbf{W} - \mathbf{W}_{r} \|_F = \sqrt{\sum_{i=r+1}^{\min(d_{\text{in}}, d_{\text{out}})} \sigma_i^2(\mathbf{W})} \le \epsilon \| \mathbf{W} \|_F
\end{equation}
که در آن $\sigma_i(\mathbf{W})$ بیانگر $i$-امین مقدار منفرد مرتب‌شده ماتریس به صورت نزولی است. نسبت حفظ انرژی طیفی $\eta_r = \frac{\sum_{i=1}^r \sigma_i^2}{\sum_{i=1}^{d} \sigma_i^2}$ تضمین می‌کند که با انتخاب آستانه $\tau \ge 0.90$، بیش از نود درصد از پراکندگی واریانس داده‌ها حفظ شده و افت سیگنال در انتشار گرادیان رو به عقب به حداقل ممکن کاهش می‌یابد.

\paragraph{تحلیل مجانبی و پیچیدگی اشتراک حافظه سخت‌افزاری}
از منظر پیچیدگی محاسباتی و انتقال داده‌ها در سلسومراتب حافظه (حافظه سراسری \lr{HBM} به حافظه محلی \lr{SRAM})، عملیات ضرب ماتریسی در این مکانیزم نیازمند تخصیص بافر موقت با ابعاد فاکتورگیری‌شده است. برای پردازش دنباله‌ای به طول $N$ با بعد مخفی $D$، بار محاسباتی هسته از فرم درجه دو به رابطه زیر تقلیل می‌یابد:
\begin{equation}
\label{eq:complexity_kronecker_tt_lora}
\mathcal{C}_{\text{kronecker\_tt\_lora}} = \mathcal{O}\left( B \cdot N \cdot D \cdot r + B \cdot H \cdot N \cdot d_k \right) \ll \mathcal{O}\left( B \cdot H \cdot N^2 \cdot d_k \right)
\end{equation}
این بهبود مرتبه پیچیدگی، با کاهش چشمگیر حجم تبادل تانسورها در گذرگاه‌های ارتباطی کارت گرافیک، مانع از اشباع پهنای باند حافظه (پدیده \lr{Memory-Bound}) در فرآیند استنتاج طولانی‌مدت می‌گردد.

\subsection{معماری و پیاده‌سازی}
ماژول در فایل \lr{Codes/framework/compression\_engine/kronecker\_tt\_lora.py} قرار داشته و شامل بخش‌های زیر است:
\begin{itemize}
    \item \textbf{تابع مقداردهی ارتوگونال و مقیاس‌گذاری طیفی (\lr{\_init\_orthogonal\_spectral}):} پارامترهای هسته تانسوری و فاکتورهای کرونکر را مستقیماً روی دستگاه هدف با استفاده از تجزیه \lr{SVD} یا \lr{QR} به صورت ارتوگونال با مقدار بیشینه تکین مقیاس‌شده به یک ($\sigma_{\max} \le 1$) مقداردهی اولیه می‌کند تا از انفجار گرادیان در آغاز آموزش جلوگیری شود.
    \item \textbf{کلاس \lr{KroneckerTTLoRALayer}:} لایه ترکیبی که به دور لایه خطی پایه تنیده شده، ماتریس پایه را منجمد نگه‌داشته و پارامترهای $\mathbf{A}_1$، $\mathbf{A}_2$، هسته‌های $\mathbf{G}_1, \mathbf{G}_2, \mathbf{G}_3$ و آداپتورهای \lr{LoRA} را آموزش می‌دهد.
    \item \textbf{سیستم کش‌گذاری برخط وزن بازسازی‌شده:} برای پرهیز از تکرار انقباض تانسوری سنگین در گذر پیشرو در طول یک تکرار ارزیابی، ماتریس $\Delta \mathbf{W}$ را به ازای ابعاد داده کش می‌کند.
    \item \textbf{کلاس \lr{KroneckerTTLoRACompressor}:} مدیریت کشف لایه‌های هدف، پیکربندی خودکار ابعاد تانسوری و اجرای بنچ‌مارک.
\end{itemize}

\subsection{نتایج تجربی}
ارزیابی ماژول \lr{kronecker\_tt\_lora} بر روی مدل‌های زبانی و انکودرهای پایه بر اساس داده‌های تجربی تاییدشده در \lr{thesis\_final/master\_verified\_results.json} انجام شده و در جدول \ref{tab:benchmark_kronecker_tt_lora} گزارش شده است:

\begin{table}[H]
\centering
\small
\caption{نتایج تجربی و اعتبارسنجی بنچ‌مارک فشرده‌ساز \lr{kronecker\_tt\_lora} روی مدل‌های پایه}
\label{tab:benchmark_kronecker_tt_lora}
\begin{tabular}{lccccc}
\toprule
\textbf{مدل پایه} \lr{(Base Model)} & \textbf{دقت} \lr{(Accuracy)} & \textbf{شاخص} \lr{(F1-Score)} & \textbf{حافظه} \lr{(Peak VRAM)} & \textbf{توان} \lr{(Throughput)} & \textbf{وضعیت} \\
\midrule
\lr{DistilBERT} & \lr{90.60\%} & \lr{0.9095} & \lr{345.95 MB} & \lr{54,737.60 tok/s} & \lr{PASS} \\
\lr{RoBERTa} & \lr{83.72\%} & \lr{0.8515} & \lr{333.23 MB} & \lr{32,763.08 tok/s} & \lr{PASS} \\
\lr{GPT-2} & \lr{55.16\%} & \lr{0.6787} & \lr{804.54 MB} & \lr{33,263.89 tok/s} & \lr{PASS\_LOW} \\
\lr{Qwen} & \lr{88.76\%} & \lr{---} & \lr{---} & \lr{---} & \lr{PASS} \\
\bottomrule
\end{tabular}
\end{table}

یافته‌های بنچ‌مارک کارآمدی ترکیبی این معماری را به خوبی نشان می‌دهد. در مدل انکودر \lr{DistilBERT} دقت بسیار بالای \lr{90.60\%} و شاخص \lr{F1} معادل \lr{0.9095} حاصل شده است. به همین ترتیب، در مدل عمیق‌تر \lr{RoBERTa} دقت عالی \lr{83.72\%} و در مدل بزرگ زبانی \lr{Qwen} دقت چشمگیر \lr{88.76\%} به دست آمده است. در مدل زبانی خودبازگشتی \lr{GPT-2} با ثبت دقت \lr{55.16\%}، وضعیت \lr{PASS\_LOW} به دلیل الزامات انطباق ساختار وزن‌های علّی گزارش شده است.

\subsection{محدودیت‌ها و تحلیل شکست}
تحلیل دقیق این ماژول پیوندی حاکی از پیچیدگی‌های ساختاری زیر است:
\begin{itemize}
    \item \textbf{انبساط پارامترهای فعال در فاز انطباق:} از آنجا که ماتریس پایه منجمد شده و وزن‌های کمکی تانسوری و آداپتورها افزوده می‌شوند، تعداد کل پارامترها در حین آموزش انطباقی نسبت به مدل بدون آداپتور اندکی افزایش می‌یابد (\lr{Moderate Expansion}) که پس از اتمام آموزش نیاز به ادغام فیزیکی (\lr{Merge}) با ماتریس اصلی دارد.
    \item \textbf{پیچیدگی محاسباتی انقباض تانسور قطار تانسوری:} بازسازی تانسور $\mathcal{TT}$ از سه هسته و ضرب هادامارد آن با ضرب کرونکر نیازمند چندین مرحله ضرب انقباضی تانسوری (\lr{Tensor Contraction}) است که موجب کاهش توان پردازشی به محدوده ۳۲ تا ۵۴ هزار توکن بر ثانیه شده است.
    \item \textbf{حساسیت شدید به مقداردهی اولیه هسته‌های تانسوری:} در صورت عدم استفاده از مقیاس‌بندی طیفی دقیق، ضرب تکراری هسته‌های قطار تانسوری با رتبه‌های میانی ممکن است دچار واگرایی یا محوشدگی مقادیر بازنمایی گردد.
\end{itemize}
